%% =============================================================================
%% D_intra Version 3 -- material for the Overleaf project
%% Generated 2026-10-01 from results/samples/{A,B,C,B_k4}{,_v3}.txt
%% Requires: \usepackage{booktabs}   (amsmath for the method equations)
%% =============================================================================


%% -----------------------------------------------------------------------------
%% 1. METHOD, Section V A 1) b -- corrected versions of the two new paragraphs.
%%    Changes from the current draft: typos ("nstead", "centriod", "demention",
%%    "In this we can"); d_h stated (the experiments use d_h = 1, which is what
%%    gives 10 -> 4); how X is handled at sampling time; why X is left out of the
%%    likelihood. Replace the existing paragraphs with these.
%% -----------------------------------------------------------------------------

\paragraph{Hop-distribution embedding}
Instead of concatenating the hop-distribution values
$\mathbf{h}_i = [h_i^1, \dots, h_i^K]$ directly into the condition vector, we
first encode them with a small MLP,
\begin{equation}
  \mathbf{e}_i^h = \mathrm{MLP}_h(\mathbf{h}_i), \qquad
  \mathrm{MLP}_h : \mathbb{R}^K \rightarrow \mathbb{R}^{d_h},
\end{equation}
and combine the encoding with the other coarse conditions,
\begin{equation}
  \mathbf{c}_i = [\, n_i,\ \rho_i,\ k,\ \mathbf{e}_i^h \,].
\end{equation}
This gives a compact representation of the radial structure before it enters
the graph-level conditioning channel $y$. $\mathrm{MLP}_h$ has one hidden layer
of width 32 with a ReLU activation, and is trained end to end with the
denoising objective. With $K = 7$ and $d_h = 1$, the dimension of the condition
falls from 10 to 4.

\paragraph{Centroid node type}
To give the hop-distribution condition a reference node, we extend the dummy
node feature from $X \in \mathbb{R}^{n \times 1}$ to
$X \in \mathbb{R}^{n \times 2}$, which identifies the centroid:
\begin{equation}
  X_v =
  \begin{cases}
    [0, 1], & v \text{ is the centroid}, \\
    [1, 0], & \text{otherwise}.
  \end{cases}
\end{equation}
Node features are observed rather than generated: they are kept fixed during
diffusion, and noise is applied only to the edge representation. They are
excluded from the likelihood for the same reason. At sampling time the first
node is designated the centroid; since the denoising network is
permutation-equivariant, this choice is without loss of generality.


%% -----------------------------------------------------------------------------
%% 2. EXPERIMENTAL SETUP -- one sentence, wherever the configurations A/B/C are
%%    described.
%% -----------------------------------------------------------------------------

We evaluate the updated model (Version 3) under exactly the same protocol as the
previous version: the same datasets, train/validation/test splits,
hyperparameters, training schedule (200 epochs) and checkpoint-selection rule. Each held-out
supernode is generated once by each model, so every comparison between the two
versions is paired on identical clusters. Differences are tested with a paired
bootstrap over clusters (2{,}000 resamples; 500 for the degree distribution).


%% -----------------------------------------------------------------------------
%% 3. TABLE -- Version 3 against the previous version. table* spans both
%%    columns of the ICDE template; it floats to the top of a page.
%% -----------------------------------------------------------------------------

\begin{table*}[t]
\centering
\caption{Version 3 (centroid node type and hop-distribution embedding) against
the previous version, evaluated on identical held-out clusters in each panel.
For the upper block of each panel, closer to the Real column is better; for the
lower block, lower is better. The last column is the 95\% paired-bootstrap
confidence interval on the difference (Version 3 $-$ previous). $^\ast$: the
interval excludes zero, i.e.\ the change is statistically distinguishable from
sampling noise; whether it is an improvement depends on the direction relative
to Real. $^\dagger$: for edge-count MAE the Real column gives the size-only
baseline $m = n - 1$. Configuration B at $k = 4$ re-samples the same Version 3
model as B at $k = 7$ on the recoarsened Minnesota test split.}
\label{tab:v3}
\begin{tabular}{lrrrc}
\toprule
& Previous & Version 3 & Real & 95\% CI on difference \\
\midrule
\multicolumn{5}{l}{\emph{Configuration A: Luxembourg $\rightarrow$ Luxembourg, $k = 7$ (882 held-out clusters)}} \\
\midrule
Connected (\%)                                 &    95.0 &    98.5 &     100.0 & $[+1.9,\,+5.2]^\ast$ \\
Acyclic (\%)                                   &    79.7 &    83.2 &      88.1 & $[+0.6,\,+6.5]^\ast$ \\
Components (mean)                              &    1.05 &    1.02 &      1.00 & $[-0.05,\,-0.02]^\ast$ \\
Edges (mean)                                   &   12.45 &   12.46 &     12.42 & $[-0.04,\,+0.06]$ \\
Edge-count MAE $\downarrow$                    &    0.30 &    0.18 & 0.19$^\dagger$ & $[-0.17,\,-0.09]^\ast$ \\
\cmidrule(lr){1-5}
Degree distribution vs.\ real $\downarrow$     &   0.036 &   0.022 &         0 & $[-0.022,\,-0.005]^\ast$ \\
Laplacian spectrum vs.\ real $\downarrow$      &   0.070 &   0.070 &         0 & -- \\
\midrule
\multicolumn{5}{l}{\emph{Configuration B: Luxembourg $\rightarrow$ Minnesota, $k = 7$ (48 held-out clusters)}} \\
\midrule
Connected (\%)                                 &    83.3 &    93.8 &     100.0 & $[-2.1,\,+22.9]$ \\
Acyclic (\%)                                   &    20.8 &    29.2 &      27.1 & $[-4.2,\,+20.8]$ \\
Components (mean)                              &    1.19 &    1.06 &      1.00 & $[-0.27,\,+0.02]$ \\
Edges (mean)                                   &   29.65 &   28.40 &     29.50 & $[-1.94,\,-0.65]^\ast$ \\
Edge-count MAE $\downarrow$                    &    1.60 &    1.35 & 2.79$^\dagger$ & $[-0.88,\,+0.35]$ \\
\cmidrule(lr){1-5}
Degree distribution vs.\ real $\downarrow$     &   0.089 &   0.218 &         0 & $[+0.071,\,+0.164]^\ast$ \\
Laplacian spectrum vs.\ real $\downarrow$      &   0.054 &   0.055 &         0 & -- \\
\midrule
\multicolumn{5}{l}{\emph{Configuration C: Luxembourg + Minnesota $\rightarrow$ Minnesota, $k = 7$ (48 held-out clusters)}} \\
\midrule
Connected (\%)                                 &    93.8 &    87.5 &     100.0 & $[-18.8,\,+6.2]$ \\
Acyclic (\%)                                   &    16.7 &    16.7 &      27.1 & $[-10.4,\,+10.4]$ \\
Components (mean)                              &    1.06 &    1.12 &      1.00 & $[-0.06,\,+0.19]$ \\
Edges (mean)                                   &   29.25 &   29.15 &     29.50 & $[-0.88,\,+0.62]$ \\
Edge-count MAE $\downarrow$                    &    1.62 &    1.85 & 2.79$^\dagger$ & $[-0.23,\,+0.69]$ \\
\cmidrule(lr){1-5}
Degree distribution vs.\ real $\downarrow$     &   0.072 &   0.113 &         0 & $[-0.009,\,+0.069]$ \\
Laplacian spectrum vs.\ real $\downarrow$      &   0.055 &   0.057 &         0 & -- \\
\midrule
\multicolumn{5}{l}{\emph{Configuration B: Luxembourg $\rightarrow$ Minnesota, $k = 4$ (114 held-out clusters)}} \\
\midrule
Connected (\%)                                 &    49.1 &    57.0 &     100.0 & $[-3.5,\,+19.3]$ \\
Acyclic (\%)                                   &    41.2 &    77.2 &      56.1 & $[+23.7,\,+47.4]^\ast$ \\
Components (mean)                              &    1.80 &    1.61 &      1.00 & $[-0.42,\,+0.02]$ \\
Edges (mean)                                   &   10.96 &   10.30 &     11.47 & $[-1.04,\,-0.27]^\ast$ \\
Edge-count MAE $\downarrow$                    &    1.30 &    1.40 & 0.82$^\dagger$ & $[-0.20,\,+0.40]$ \\
\cmidrule(lr){1-5}
Degree distribution vs.\ real $\downarrow$     &   0.105 &   0.250 &         0 & $[+0.094,\,+0.177]^\ast$ \\
Laplacian spectrum vs.\ real $\downarrow$      &   0.117 &   0.105 &         0 & -- \\
\bottomrule
\end{tabular}
\end{table*}


%% -----------------------------------------------------------------------------
%% 4. TABLE -- do generated clusters follow the hop distribution? Fits one
%%    column.
%% -----------------------------------------------------------------------------

\begin{table}[t]
\centering
\caption{Fidelity to the hop-distribution condition: total-variation distance
between the hop distribution of the generated cluster, measured from its
centroid, and the requested $\mathbf{h}_i$ (0 is exact), and the share of
clusters that match exactly. For Version 3 the distance is measured from the
designated centroid. The previous version has no designated centroid, so it is
credited with the best-matching node of each generated cluster, which favours
it.}
\label{tab:v3-hop}
\begin{tabular}{lrrrr}
\toprule
& \multicolumn{2}{c}{Previous (best node)} & \multicolumn{2}{c}{Version 3 (centroid)} \\
\cmidrule(lr){2-3}\cmidrule(lr){4-5}
& TV $\downarrow$ & Exact (\%) & TV $\downarrow$ & Exact (\%) \\
\midrule
A, $k = 7$ & 0.060 & 52.4 & 0.085 & 35.1 \\
B, $k = 7$ & 0.106 & 10.4 & 0.164 &  6.2 \\
C, $k = 7$ & 0.122 &  4.2 & 0.186 &  4.2 \\
B, $k = 4$ & 0.154 & 14.9 & 0.277 &  5.3 \\
\bottomrule
\end{tabular}
\end{table}


%% -----------------------------------------------------------------------------
%% 5. RESULTS TEXT
%% -----------------------------------------------------------------------------

\paragraph{Effect of Version 3}
Table~\ref{tab:v3} compares Version 3 with the previous version. In
distribution (configuration A, 882 held-out Luxembourg clusters), Version 3 is
significantly better on every structural measure except the Laplacian spectrum,
which is unchanged. Connectivity rises from 95.0\% to 98.5\% and acyclicity
from 79.7\% to 83.2\% (real: 88.1\%); the edge-count error falls from 0.30 to
0.18, so the model now beats the size-only baseline $m = n - 1$ (0.19); and the
degree-distribution distance falls from 0.036 to 0.022. Adding the hop
distribution to the condition had previously cost connectivity and acyclicity
relative to conditioning on $[n_i, \rho_i, k]$ alone; Version 3 recovers most
of that loss while retaining the hop condition.

Under transfer to Minnesota the effect is mixed. At $k = 7$ (48 clusters) the
changes in connectivity and acyclicity for B are not significant, but Version 3
generates significantly fewer edges (28.40 against 29.65; real 29.50) and its
degree-distribution distance increases from 0.089 to 0.218. No difference is
significant for C. At $k = 4$ (114 clusters) the same model again generates
too few edges (10.30 against 10.96; real 11.47), and acyclicity overshoots the
real value, rising from 41.2\% to 77.2\% against 56.1\%. Version 3 therefore
improves in-distribution generation but biases transferred clusters towards
sparser, more tree-like structures.

Table~\ref{tab:v3-hop} shows that Version 3 does not reproduce the requested
hop distribution more closely than the previous version: measured from the
designated centroid, 35.1\% of in-distribution clusters match it exactly,
against 52.4\% for the previous version credited with its best node. A
plausible explanation is that $\mathrm{MLP}_h$ with $d_h = 1$ compresses the
six degrees of freedom of $\mathbf{h}_i$ into a single scalar, so the model
receives only a summary of the radial profile. Separating the two changes,
the centroid node type with the full $\mathbf{h}_i$, is left to an ablation.
